\section{The persistence problem}

The motivating question is elementary to state and difficult to answer without
slipping between meanings of information: how can the consequence of an
interaction persist after the object or arrangement that first carried it has
changed or disappeared? A footprint can outlast a foot; a measurement record
can be copied into an environment; angular momentum can be transferred among
spin, orbital, and field degrees of freedom; a biological response can survive
component turnover. None of these examples requires information to float free
of physical realization. They instead distinguish persistence beyond an
\emph{original} carrier from existence without \emph{any} carrier.

IOTA takes that distinction as a research problem rather than a slogan. Its
organizing hypothesis is that information is the interaction-caused change in
the reachable. In condition form, the information present at time $t$ is the
part of realized interaction history that remains causally active in present
local state and relational topology. In change form, an interaction is
informational when it changes which future trajectories or responses are
reachable, or changes their probabilities. The proposal is intentionally
stronger than Shannon information as uncertainty and narrower than the claim
that every physical difference is information.

This paper gives the synthesis two jobs. The first is diagnostic: does this
program have genuine scientific legs when compared with the best existing
science, including theories that may already contain its central object? The
second is constructive: if it does, what sequence of results would move it from
a bounded model of persistent constraint toward an account of familiar physics?
The review is therefore allowed to return a negative verdict. A literature
collection selected only because it sounds information-friendly would defeat
both purposes.

\subsection{IOTA as substrate hypothesis, not a frozen cosmic lattice}

The IOTA substrate hypothesis treats informational events or distinctions,
relations, and local rules of lawful succession as the irreducible explanatory
layer. The concrete lattice in the present simulation is a candidate
realization, not a claim that its dimensionality, neighborhood, scheduler, or
constants are the uniquely correct microscopic architecture of nature. This
distinction is essential. A literal fixed lattice inherits immediate burdens:
preferred structure, continuum and symmetry recovery, dependence on update
order, and the possibility that every reported coherence is a special feature
of one engineered rule.

At the theoretical level, write the retained present condition as
\begin{equation}
 H_{\le t}\longrightarrow M_t=(X_t,G_t),
 \label{eq:history-record}
\end{equation}
where $H_{\le t}$ is realized interaction history, $X_t$ the current local
state, and $G_t$ the active relational organization. Fixed law, topology,
scheduler, boundary, constants, and exogenous inputs must be declared rather
than hidden in the word ``system.'' Two histories are operationally equivalent
at $t$ only when no allowed future intervention distinguishes the retained
states at the chosen scale and precision.

The ontology claim is a postulate under investigation: familiar physical
objects and quantities are proposed to be stable effective organizations of
this underlying structure. It is not inferred from the mere use of a computer
model, nor from the truism that physical theories can be encoded as data.

\subsection{ICP and the emergence ladder}

ICP is the proposed principle connecting the substrate hypothesis to dynamics:

\begin{quote}
\textbf{Physically effective information is interaction-formed, embodied
constraint on an explicitly indexed family of reachable or probable futures;
it persists when that constraint role is lawfully re-instantiated beyond its
original carrier.}
\end{quote}

The proposal has a definite explanatory order:
\begin{enumerate}
  \item local distinctions and relations make some successor events lawful and
  others unavailable;
  \item interaction changes that successor structure;
  \item some organizations preserve and transmit selected constraints;
  \item carrier succession allows a relational or causal role to outlast its
  original realization;
  \item persistent coherences may acquire stable collective invariants;
  \item suitable coherences may realize quasiparticle and field roles; and
  \item mass, energy, geometry, clocks, and macroscopic physics may appear as
  effective properties of those organizations.
\end{enumerate}

Only the beginning of this ladder has project evidence. The current quantum
Dirac construction starts higher up: it assumes quantum amplitudes,
composition, a Born readout, a clock, and a mass schedule, then tests whether
specified local updates recover a bounded known role
\citep{Clark2026DiracRole}. It is evidence that a local construction can realize
that role under supplied quantum structure, not yet a derivation of quantum
quasiparticles from the base IOTA lattice. The gap is not a footnote. It is the
central research frontier.

\subsection{Why the synthesis must be adversarial}

The phrase ``constrains the future'' is dangerously easy to satisfy. Given a
complete physical state and law, virtually every state difference can alter a
future trajectory. Control theory already studies reachable sets; statistical
models already condition future distributions; causal models already compare
interventions; and semantic-information work already erases correlations and
measures effects on later viability. If ICP merely renames these objects, it may
be useful language but is not a new physical principle.

The burden is therefore not to find fields in which constraints matter. It is
to determine whether the combined requirements of historical formation,
physical embodiment, sufficiency for a declared intervention-response family,
causal carrier manipulation, succession across original-carrier replacement,
and representation robustness produce a new invariant, compression,
prohibition, or prediction.
